\subsection{A third problem and a synthesis across environments}
\label{sec:pfsp}

To test whether the attribution is a property of parallel machine scheduling, the protocol was applied to the permutation flowshop (makespan), the problem in which Iterated Greedy originated \citep{RuizStutzle2007}.
The IG starts from NEH, removes $d$ random jobs, reinserts them at the best position (Taillard's acceleration), applies an insertion local search and accepts with a constant temperature; the action is the
destruction size $d\in\{2,3,4,6,8\}$. The ILS has the same local search and perturbs by two random insertions. The instances are Taillard-style, generated here ($n\in\{50,100\}$, $m\in\{10,20\}$; $30$ test and $10$ tuning
instances per size; $T=300$); comparisons are paired, CPU times are equalised by sequential calibration, and the design, the Holm family of five tests and the predictions were recorded before the test (Table~\ref{tab:f36_pfsp}).

\paragraph{The active ingredient is yet another}
In the flowshop the destruction--reconstruction cycle does not separate the methods: the IG and the ILS, with the same local search, are equivalent at equal CPU time ($-0.03\%$, $p_{\mathrm{Holm}}=0.18$; the
prediction that the IG would win was wrong). The local search, however, pays its cost: the IG without it, with $17$ to $57$ times more iterations for the same CPU time, is worse by $0.21\%$ ($65$ wins against $17$,
$p_{\mathrm{Holm}}=10^{-6}$), the opposite of what was found for identical machines. The active ingredient is thus the destruction--reconstruction cycle for identical machines, its combination with a descent for unrelated
machines and the local search for the flowshop, and only an ablation at equal CPU time reveals which.

\paragraph{The policy over the destruction size is harmful again}
Q-Learning over $d$ is worse than random choice of $d$ ($+0.14\%$, $p_{\mathrm{Holm}}=5\times10^{-4}$), the validated size beats it by $0.16\%$ ($p_{\mathrm{Holm}}=8\times10^{-5}$), and the headroom of the size over
random choice is null ($-0.02\%$). The signs of the predictions recorded beforehand were right, except that the validated size was expected only to match the agent.

% Auto-generated by analysis/make_tables_cycle5.py (Entrada 129)
\begin{table}[t]
\centering
\caption{A third problem: the permutation flowshop (Taillard-style instances generated here, $n\in\{50,100\}$, $m\in\{10,20\}$; 30 test instances per size, 10 tuning instances per size; Iterated Greedy of Ruiz and St\"utzle with destruction size $d\in\{2,3,4,6,8\}$ as the action; $T=300$). $\bar d$ in \% of the makespan (negative: first better); Wilcoxon, Holm over five. The ILS has the same insertion local search; the IG without local search receives $T_{\mathrm{eq}}$ iterations (5{,}177 to 17{,}076) so that the CPU time is the same, and the ILS $217$ to $250$.}
\label{tab:f36_pfsp}
{\scriptsize\setlength{\tabcolsep}{3.5pt}
\begin{tabular}{@{}lrrrr@{}}
\toprule
Comparison & W/L/T & $\bar d$ (\%) & 90\% interval & $p_{\mathrm{Holm}}$ \\
\midrule
P1: IG vs.\ ILS at equal CPU (mechanism) & 53/31/36 & -0.027 & [-0.070, +0.014] & 0.185 \\
P2: IG with vs.\ without local search at equal CPU & 65/17/38 & -0.212 & [-0.279, -0.148] & $1\times10^{-6}$ \\
P3: \textsc{QL} over $d$ vs.\ random $d$ (equal $T$) & 25/68/27 & +0.142 & [+0.073, +0.210] & $5\times10^{-4}$ \\
P4: validated $d$ vs.\ \textsc{QL} & 64/25/31 & -0.161 & [-0.225, -0.101] & $8\times10^{-5}$ \\
P5: validated $d$ vs.\ random $d$ & 47/38/35 & -0.020 & [-0.083, +0.043] & 0.466 \\
\bottomrule
\end{tabular}}
\end{table}


\paragraph{Replication on the classic instances}
The study was repeated on $40$ classic instances of \citet{Taillard1993} (ta041--ta050, ta051--ta060, ta071--ta080, ta081--ta090; regenerated from the published generator and seeds and checked against the matrix printed with it),
with every setting frozen from the study above (the validated size, state thresholds and CPU ratios, all chosen on generated instances; realised CPU ratios $1.06$ and $1.09$), the same family of five tests and the predictions recorded beforehand.
The findings replicate. The IG and the ILS are equivalent at equal CPU time ($-0.006\%$, $p_{\mathrm{Holm}}=0.74$), the local search pays its cost ($-0.28\%$, $26$ wins against $7$, $p_{\mathrm{Holm}}=7\times10^{-4}$), Q-Learning over $d$ is not better than random choice
($+0.065\%$, interval $-0.02$ to $+0.16$), the validated size beats the agent by $0.17\%$ (interval $-0.29$ to $-0.05$; $p_{\mathrm{Holm}}=0.064$, short of the threshold with $40$ instances) and its headroom over random choice is $-0.11\%$ ($p_{\mathrm{Holm}}=0.12$; the magnitude predicted beforehand was below $0.05\%$). With $40$ instances only the local-search effect is individually confirmed, so the replication is of signs and magnitudes (Supplementary Section~\ref{S-sec:m-taillard}).

\paragraph{Synthesis}
Table~\ref{tab:f35_synthesis} puts the seven environments side by side. The oracle headroom is negative (some fixed action beats random selection per instance) in six of the seven, and the validated choice
captures a part of it wherever it is real. The agent behaves differently by type of action. Over qualitatively different destruction \emph{strategies} its gain is small but never positive, and it recovers
$0.44$ to $0.68$ of the validated gain at $1000\times50$ and on the standard benchmark. Over the destruction \emph{size}, an ordered parameter, it is worse than random choice in all three problems
($+0.015\%$, $+0.152\%$, $+0.142\%$), although the size has headroom that validation captures in two of them. This regularity was noticed after the results and is a hypothesis, not a finding: headroom is necessary for a
policy gain, and whether an online agent can use it seems to depend on the type of action and on how early the reward arrives.

% Auto-generated by analysis/make_table_synthesis.py (Entrada 154)
\begin{table}[t]
\centering
\caption{Headroom and policy gain across the environments of the study (mean $\bar d$ in \% of the makespan relative to random selection with the same budget; negative: better). Oracle: best of the fixed actions per instance (an upper bound that includes selection noise); validated: the action chosen offline on separate instances; Q-Learning: the online agent (tabular, $T=1000$ except the flowshop, $T=300$). The last column is the CPU time of Q-Learning relative to random selection.}
\label{tab:f35_synthesis}
{\scriptsize\setlength{\tabcolsep}{3pt}
\begin{tabular}{@{}lrrrrrr@{}}
\toprule
Environment & $N$ & Actions & Oracle & Validated & Q-Learning & CPU ratio \\
\midrule
$\mathrm{P}||C_{\max}$, destruction strategy & 1{,}130 & 3 & -0.039 & +0.001 & -0.016 & -- \\
$\mathrm{P}||C_{\max}$, destruction size & 1{,}200 & 5 & -0.127 & -0.089 & +0.015 & -- \\
$\mathrm{R}||C_{\max}$ synthetic, destruction strategy & 400 & 4 & -0.240 & -- & -0.073 & 0.97 \\
$\mathrm{R}||C_{\max}$ $1000{\times}50$, destruction strategy & 96 & 4 & +0.092 & -0.461 & -0.312 & 1.42 \\
$\mathrm{R}||C_{\max}$ standard benchmark, destruction strategy & 700 & 4 & -0.149 & -0.124 & -0.055 & 1.12 \\
$\mathrm{R}||C_{\max}$ standard benchmark, destruction size & 700 & 5 & -0.325 & -0.201 & +0.152 & 0.65 \\
Permutation flowshop, destruction size & 120 & 5 & -0.193 & -0.020 & +0.142 & 0.89 \\
\bottomrule
\end{tabular}}
\end{table}


\paragraph{The learners of other classes in the other two problems}
The contextual bandit and the neural Q-learner of Section~\ref{sec:unrelated} were also run in the other two problems, over the destruction size, with the same instances, actions and baselines (Table~\ref{tab:f40_online}; predictions recorded beforehand). For identical machines neither differs from random choice
($+0.003\%$ and $+0.006\%$; $p_{\mathrm{Holm}}=0.82$) and validation beats both by $0.09\%$ ($p_{\mathrm{Holm}}<10^{-30}$). In the flowshop LinUCB is level with random choice ($+0.027\%$), the neural learner is significantly \emph{worse} ($+0.099\%$, $p_{\mathrm{Holm}}=0.027$) and validation
beats it by $0.14\%$ ($p_{\mathrm{Holm}}=4\times10^{-4}$) and LinUCB by $0.07\%$ ($p_{\mathrm{Holm}}=0.15$, not confirmed with $160$ instances). The attribution does not change in any of the three problems.

% Auto-generated by analysis/make_table_online.py (Entrada 167/168)
\begin{table}[t]
\centering
\caption{The contextual bandit (LinUCB) and the neural Q-learner choosing the destruction size in the two other problems (same instances, seeds, actions and baselines as before; hyperparameters chosen on tuning instances of each problem; design and predictions recorded beforehand). $\bar d$ in \% of the makespan (negative: first better); Wilcoxon, Holm over four within each problem; confirmed: $p_{\mathrm{Holm}}<0.05$ and $\bar d<0$ (the neural learner is significantly \emph{worse} than random choice in the flowshop: $\bar d>0$).}
\label{tab:f40_online}
{\scriptsize\setlength{\tabcolsep}{3pt}
\begin{tabular}{@{}lrrrrl@{}}
\toprule
Comparison & W/L/T & $\bar d$ (\%) & 90\% interval & $p_{\mathrm{Holm}}$ & Confirmed \\
\midrule
\multicolumn{6}{l}{P$||C_{\max}$ ($1{,}200$ instances, $T=1000$)} \\
\quad LinUCB vs random $d$ & 161/170/869 & +0.003 & [-0.008, +0.014] & 0.82 & no \\
\quad neural Q-learning vs random $d$ & 163/174/863 & +0.006 & [-0.005, +0.017] & 0.82 & no \\
\quad validated $d$ vs LinUCB & 338/53/809 & -0.091 & [-0.102, -0.080] & $4\times10^{-31}$ & yes \\
\quad validated $d$ vs neural Q-learning & 334/44/822 & -0.094 & [-0.106, -0.083] & $2\times10^{-33}$ & yes \\
\midrule
\multicolumn{6}{l}{Flowshop ($160$ instances: $120$ generated, $40$ Taillard; $T=300$)} \\
\quad LinUCB vs random $d$ & 58/58/44 & +0.027 & [-0.021, +0.073] & 0.44 & no \\
\quad neural Q-learning vs random $d$ & 48/72/40 & +0.099 & [+0.046, +0.153] & 0.03 & no \\
\quad validated $d$ vs LinUCB & 65/57/38 & -0.069 & [-0.133, -0.009] & 0.15 & no \\
\quad validated $d$ vs neural Q-learning & 79/38/43 & -0.141 & [-0.198, -0.084] & $4\times10^{-4}$ & yes \\
\bottomrule
\end{tabular}}
\end{table}

